% chapter3-beamer.tex — a standalone slide deck for Chapter 3.
% See chapter1/beamer/chapter1-beamer.tex for the full explanation of the
% `slim` option and the adapt-don't-\input approach.  Build with
% `make chapter3-beamer`.
%
% Chapter 3's prose is mostly placeholder, so this deck is a skeleton: a title,
% an overview, one result frame, and a discussion frame.  Replace the contents
% with your actual results — the notation macros are already available.

\documentclass{beamer}

\usetheme{Madrid}
\setbeamertemplate{navigation symbols}{}
\setbeamertemplate{theorems}[numbered]

\usepackage[slim]{philogic}

\usepackage[backend=biber]{biblatex}
\addbibresource{../../refs.bib}

\title{Main Results}
\subtitle{Chapter 3}
\author{Your Name}
\institute{Your Department \\ University of California, Davis}
\date{}

\begin{document}

\frame{\titlepage}

\begin{frame}{Overview}
  \begin{itemize}
    \item A first result, building on the framework of Chapters 1--2.
    \item A second result, extending the analysis.
    \item Discussion and consequences.
  \end{itemize}
\end{frame}

\begin{frame}{First result}
  Building on the relational models \MRel{} of Chapter~1 and the neighbourhood
  models \MN{} of Chapter~2, we show that\dots
  \begin{theorem}
    \emph{(State your central result here.)}  For every formula $A$,
    $\proves A$ if and only if $\models A$.
  \end{theorem}
\end{frame}

\begin{frame}{Discussion}
  \begin{itemize}
    \item What the result buys us, and how it bears on \K, \KT, and \Sfour.
    \item Connections to neighbourhood semantics for non-normal logics
      \parencite{pacuit2017}.
    \item Open questions, taken up in the conclusion.
  \end{itemize}
\end{frame}

\begin{frame}[allowframebreaks]{References}
  \printbibliography
\end{frame}

\end{document}
